\section{Passage to Euclidean space and Sobolev estimates}\label{sec:transfer}

The aperture in Proposition~\ref{prop:periodic-cap} is fixed. A finite
partition of the sphere therefore gives the critical estimate on the
whole annulus without increasing its logarithmic power. We first remove
the spatial period and then account for the cost of summing annuli.

\subsection{Removing the period}

\begin{proof}[Proof of Theorem~\ref{thm:main}]
Write $h=N^{-1}$ and $P_h=\varphi(hD)$, with $\varphi$ supported
in a fixed compact annulus. Choose a smooth finite partition
$\varphi=\sum_\nu\varphi_\nu$ subordinate to the admissible caps.
The Euclidean kernel of $\varphi_\nu(hD)$ is
$h^{-n}\check\varphi_\nu(\cdot/h)$, and periodization does not increase
its $L^1$ norm. Apply Proposition~\ref{prop:periodic-cap} to
$\varphi_\nu(hD)f$ on the same torus for every $\nu$, and sum. This gives
\begin{equation}\label{eq:periodic-critical}
 \|\U(T)P_h f\|_{L^{\pc}(I\times\TT^n)}
 \lesssim [\log(2/h)]^{2/\pc}\|f\|_{L^{\pc}(\TT^n)},\qquad 0<h\le1.
\end{equation}

The Euclidean convolution kernel of $\U(T)P_h$ is
\[
 K_h(T,X)=(2\pi h)^{-n}
 \int e^{\mathrm i(X\cdot\eta-T|\eta|)/h}\varphi(\eta)\dd\eta .
\]
For $T\in I$ and $|X|\ge3$ the gradient of the phase before division
by $h$ has length at least $|X|-2\ge|X|/3$. Integration by parts in
the gradient direction gains $h/|X|$ each time: the amplitude is
supported away from zero and all phase derivatives of order at least
two are bounded there. Using at least $M+n$ integrations gives
\begin{equation}\label{eq:off-cone-kernel}
 |K_h(T,X)|\le C_M h^M(1+|X|)^{-M},
 \qquad T\in I,\quad |X|\ge3 .
\end{equation}

Tile $\RR^n$ by unit cubes $Q$, and write $10Q$ for the concentric
cube of side ten. On $Q$ the contribution of $f\mathbf1_{(10Q)^c}$
is bounded pointwise by
$C_M h^M ((1+|\cdot|)^{-M}*|f|)(X)$, since $|X-Y|\ge9/2$ for
$X\in Q$ and $Y\notin10Q$. This single function bounds the far
contributions on every cube. For $M>n$, its $L^{\pc}(I\times\RR^n)$
norm is $O(h^M)\|f\|_{\pc}$ by Young's inequality.

Translate the center of $Q$ to the origin and periodize
$f_Q=f\mathbf1_{10Q}$ on the torus of side $100$. Its torus norm is
$\|f\|_{L^{\pc}(10Q)}$. The periodic wave kernel is the periodization
of $K_h$. On $Q$, its nonzero translates see only arguments
$X-Y+100k$ with $Y\in10Q$ and $k\in\mathbb Z^n\setminus\{0\}$.
Since $|X_j-Y_j|\le5.5$, these arguments have length at least
$94.5|k|$. Summing \eqref{eq:off-cone-kernel} over $k$ bounds the
difference between the two outputs there by $C_M h^M\|f_Q\|_1$.
As $10Q$ has fixed volume, \eqref{eq:periodic-critical} now gives
\[
 \|\U(T)P_h f_Q\|_{L^{\pc}(I\times Q)}
 \lesssim [\log(2/h)]^{2/\pc}\|f\|_{L^{\pc}(10Q)} .
\]
Take the $\pc$-th powers and sum over $Q$. The cubes $10Q$ have
bounded overlap. Combining this sum with the far contributions gives
the asserted Euclidean estimate, since $\log(2/h)\asymp\log(2+N)$
for $N\ge2$.
\end{proof}

\subsection{Interpolation and summation}

Plancherel bounds the norm of $\U(T)P_h$ from $L^2_X$ to $L^2_{T,X}$
by $\|\varphi\|_\infty$, since $|I|=1$. At $p=\infty$ its norm is
$O(h^{-d/2})$. To see this, partition the annulus into $O(h^{-d/2})$
caps of aperture $c\sqrt h$ with bounded overlap. In coordinates
$\eta=(\sqrt h\,z,\rho)$ about the axis of a cap, its cutoff is
supported in a fixed compact set with $\rho$ bounded below and has
uniformly bounded derivatives. Such cutoffs are obtained by dividing
rescaled angular bumps by their sum; the overlapping centers have
bounded offsets in these coordinates. Removing the linear phase
$T\rho/h$ leaves
\[
 \frac{T}{h}\bigl(\sqrt{\rho^2+h|z|^2}-\rho\bigr)
 =\frac{T|z|^2}{\sqrt{\rho^2+h|z|^2}+\rho}.
\]
The expression on the right has uniformly bounded derivatives on the
support. Integration by parts thus bounds the $L^1$ norm of each
rescaled kernel by a constant. The affine change of frequencies and
the spatial translation from the linear phase preserve that norm.
Summing over the caps gives the claimed bound.

Set $s_+(p)=\max\{0,d/2-n/p\}$ and
$\beta(p)=\min\{d(1-2/p),2/p\}$. Interpolating
Theorem~\ref{thm:main} with the two bounds above gives
\begin{equation}\label{eq:all-p-frequency}
 \|\U(T)P_h f\|_{L^p(I\times\RR^n)}
 \lesssim_p h^{-s_+(p)}[\log(2/h)]^{\beta(p)}\|f\|_{L^p(\RR^n)},
 \qquad 2\le p\le\infty .
\end{equation}
Indeed, above $\pc$ the power of $h^{-1}$ is
$(d/2)(1-\pc/p)=d/2-n/p$. We use Riesz--Thorin interpolation here
and the spatial Littlewood--Paley and Mihlin theorems below; see
\cite[Chapters~1 and~6]{Grafakos2014}.

\begin{proof}[Proof of Theorem~\ref{thm:local-smoothing}]
The gain of an arbitrary $\varepsilon>0$ now allows an absolute
sum over dyadic annuli. Indeed, set $g=J^{s_+(p)+\varepsilon}f$.
At frequency $2^j$, $j\ge1$, \eqref{eq:all-p-frequency} bounds the
output by $C2^{-j\varepsilon}(1+j)^{\beta(p)}\|g\|_p$: a smooth
enlargement of the annulus times $2^{j(s_+(p)+\varepsilon)}
J^{-s_+(p)-\varepsilon}$ has uniformly integrable kernel.
The coefficients are summable. On the remaining low frequencies the
cutoff times $e^{-\mathrm iT|\xi|}$ obeys uniform Mihlin bounds for
$T\in I$. At $p=2$ use Plancherel without any loss.
\end{proof}

The same summation can be made more precise if one wants an estimate
at the exact power $s_+(p)$.

\begin{corollary}[Logarithmic Sobolev estimate]\label{cor:log-sobolev}
Let $2<p<\infty$, $\delta_p=1/2-1/p$, and $B_*(p)=\beta(p)+\delta_p$.
Write $J=(1-\Delta)^{1/2}$, $\ell_1(r)=\log(2+r)$ and
$\ell_2(r)=\log(2+\ell_1(r))$. Then
\begin{equation}\label{eq:log-sobolev-summed}
 \|\U(T) f\|_{L^p(I\times\RR^n)}
 \lesssim_{p,B,C} \|J^{s_+(p)}\ell_1(J)^B\ell_2(J)^C f\|_{L^p(\RR^n)}
\end{equation}
whenever $B>B_*(p)$ and $C\in\RR$, or $B=B_*(p)$ and $C>\delta_p$.
At $p=\pc$, $B_*(\pc)=1/2+1/\pc$ and $\delta_{\pc}=1/(2n)$.
At $p=2$ no logarithmic factor is needed.
\end{corollary}

\begin{proof}[Proof of Corollary~\ref{cor:log-sobolev}]
Let $P_j$ be the projections of a smooth dyadic partition of the high
frequencies, $|\xi|\asymp 2^j$, $j\ge1$, and $\widetilde P_j$ an
enlarged cutoff equal to one on the support of the symbol of $P_j$.
Write $s=s_+(p)$ and
$g=J^s\ell_1(J)^B\ell_2(J)^C f$. On the $j$-th annulus put
\[
 A_j=2^{js}(1+j)^B[\log(2+j)]^C\,
      \widetilde P_j J^{-s}\ell_1(J)^{-B}\ell_2(J)^{-C}.
\]
If $a_j$ is the symbol of $A_j$, then the functions $a_j(2^j\eta)$
have common compact annular support and uniformly bounded derivatives.
Bounded overlap
of the original annuli therefore gives uniform Mihlin bounds for
all signed sums of the $A_j$. Randomizing these signs and using
Khintchine's inequality gives, for $p\ge2$,
\[
 \Big(\sum_j\|A_jg\|_p^p\Big)^{1/p}
 \le \Big\|\Big(\sum_j |A_jg|^2\Big)^{1/2}\Big\|_p
 \lesssim_p \|g\|_p .
\]

The waves $u_j(T)=\U(T)P_j f$ retain their spatial frequency supports.
The reverse Littlewood--Paley inequality at each $T$, followed by the
triangle inequality in $L^{p/2}(I\times\RR^n)$, gives
\[
 \Big\|\sum_j u_j\Big\|_{L^p_{T,X}}
 \lesssim_p \Big(\sum_j\|u_j\|_{L^p_{T,X}}^2\Big)^{1/2}.
\]
By \eqref{eq:all-p-frequency}, the $j$-th norm on the right is at most
$C w_j\|A_jg\|_p$, where
$w_j=(1+j)^{\beta(p)-B}[\log(2+j)]^{-C}$.
With $1/r=\delta_p$, H\"older's inequality bounds this sequence in
$\ell^2$ by $\|w\|_{\ell^r}(\sum_j\|A_jg\|_p^p)^{1/p}$. The weight
belongs to $\ell^r$ when $B>B_*(p)$. At $B=B_*(p)$ it belongs to
$\ell^r$ precisely when $\sum_{j\ge2} j^{-1}(\log j)^{-Cr}<\infty$,
that is, $C>\delta_p$.

For the remaining low-frequency cutoff $\chi_0$, the symbol
$e^{-\mathrm iT|\xi|}\chi_0(\xi)$ obeys the Mihlin derivative bounds
uniformly in $T\in I$ away from $\xi=0$. The same holds after
multiplication by the inverse of the Sobolev and logarithmic weights,
which are smooth and positive near zero. This proves
\eqref{eq:log-sobolev-summed}, first for finite frequency sums and
then by density. At $p=2$ use conservation of the $L^2$ norm.
\end{proof}

The condition on $B$ and $C$ records the cost of this frequency
summation. It need not be optimal for the wave operator: for
$n\ge4$ and $p>2+4/(n-3)$, Heo, Nazarov and Seeger proved
\eqref{eq:log-sobolev-summed} with $B=C=0$
\cite[Theorem~1.2]{HeoNazarovSeeger2011}.
